% ЕМТ 2008, VIII клас — проверен LaTeX препис на липсващата задача 8.1 от официалните файлове в Klasirane.com.
% Условия: https://klasirane.com/api/competitions/EMT/All/2008/8%20%D0%BA%D0%BB/probs
% SHA-256 (условия): e2b69d5a3b8fa5bf6aabf2cecf11a963acbcaa8283248dd4f188e9bfccbf91c0
% Решения: https://klasirane.com/api/competitions/EMT/All/2008/8%20%D0%BA%D0%BB/sol
% SHA-256 (решения): 1022da9f3e90121ba046b644e9fd00c94225c2a1976d6959e6167ef220d65163
% Mathpix Snip (условия): 6ca71a09-9cfa-4771-a9cf-7464c6a0b44a
% Mathpix Snip (решения): 20fc0f70-6e6e-4527-be0d-5c386e8411cf
% Формулировката и решението са сверени визуално с официалните PDF страници; задачата няма фигура.

Задача 8.1. Да се реши уравнението $|x-m|+|x+m|=x$ в зависимост от стойностите на параметъра $m$.

Решение. Нека $x$ е едно решение на уравнението. Тогава от неравенството за модулите ще получим $x=|x-m|+|x+m| \geq|(x-m)+(x+m)|=2|x|$, което е възможно само при $x=0$. Но при $x=0$ ще имаме $m=0$. Окончателно уравнението има решение $x=0$ при $m=0$ и няма решение при $m \neq 0$.
